\scanpage{184}\mspage{39}%
% archive sheet number: 184
% L1
on arrive à \unsure{réconstituer|rêconstituer} le type d'homotopie
% L2
via son $\infty$-groupoïde fondamental, comme
% L3: the doubtful older loop lies underneath the beginning of live et.
expliqué dans ma lettre précédente, \deleted{\unsure{e|a}}et
% L4
rappelé dans celle-ci. (Cela signifie en
% L5--L6: physical line-end coho- is joined.
particulier que toutes les opérations coho%
mologiques possibles et imaginables sont
% L7
déjà incluses dans la donnée d'un
% L8
tel système de $n$-catégories .\,.\,.\,.) \deleted{Mais}

% L9
De même, les types d'homotopie plus
% L10
\unsure{cannelés|canelés}, liés à certains diagrammes finis,
% L11
qu'on peut associer à \deleted{\unsure{des}}certains types
% L12/L12a: P is raised in the gap after stratifications; no separate caret.
de stratifications \inserted{$\unsure{\mathcal P|P}$} d'espaces topologiques
% L13: X shares the main baseline. A small separate point at (702,1186)
% is above the latter part of disons; its origin and function are unassigned.
``modérés'' $X$ disons\raisebox{0pt}[0pt][0pt]{\parbox{0pt}{\figure{\begin{tikzpicture}[overlay]\fill (-0.85em,1.35em) circle[radius=0.32pt];\end{tikzpicture}}}}, devraient correspondre de
% L14/L14a: (n+1)- replaces the compressed baseline prefix.
% Separate point (1156,1367), below the g-descender of catégories.
façon aussi parfaite aux \deleted{\unsure{$n$-}}\inserted{$(n+1)$-}catég\raisebox{0pt}[0pt][0pt]{\parbox{0pt}{\figure{\begin{tikzpicture}[overlay]\fill (0,-0.85em) circle[radius=0.4pt];\end{tikzpicture}}}}ories de $n$-champs
% L15
sur $X$ qui sont loc constants sur chacune des
% L16: the upper point (1622,1482) is physically separate from the
% closing parenthesis; the lower point after that parenthesis is literal.
strates de $\mathcal P$ (disons: qui sont \unsure{subordonnés} à $\mathcal P$\raisebox{0pt}[0pt][0pt]{\parbox{0pt}{\figure{\begin{tikzpicture}[overlay]\fill (0.5em,1em) circle[radius=0.3pt];\end{tikzpicture}}}}).
% L17
Si plus haut le \unsure{topis|topos} de la cat. dérivée loc.
% L18--L19: physical line-end s'atten- is joined.
constante marchait bel et bien, on s'atten%
drait ici à récupérer le type d'homotopie
% L20/L20a: ss- has a curved guide to the beginning of catégorie.
% Two separate points (1001,1963), (996,1970) lie below and slightly left
% of that beginning; they are not a colon or a mathematical operator.
mixte à partir de la \inserted{\unsure{ss-|s-}}\raisebox{0pt}[0pt][0pt]{\parbox{0pt}{\figure{\begin{tikzpicture}[overlay]\fill (-0.15em,-1.1em) circle[radius=0.32pt];\fill (-0.3em,-1.28em) circle[radius=0.32pt];\end{tikzpicture}}}}catégorie de la \unsure{catégori|catégorie}
% L21
dérivée des faisceaux abéliens sur $X$,
% L22: the reduced two-component sign is preferably à; colon shape retained.
formée des complexes qui sont \unsure{à|:} ch. loc.
% L23
\unsure{constant|constante} sur chacune des strates -- avec
% L24: underline starts just before H and ends just beyond m; R is outside.
en plus les opérations $\overset{\mathbb{L}}{\otimes}$, $\mathbb{R}\underline{\mathrm{Hom}}$ plus, \unsure{les|le}

